\chapter{Conclusion and Future Work}

\section{Conclusion and Implications}
This thesis investigates chest X-ray diagnostic classification using two standard convolutional neural network models, ResNet50 and DenseNet121, under both scratch-training and transfer-learning settings. The study is designed as a focused and reproducible comparison on NIH ChestXray14 using a binary normal-versus-abnormal formulation. In the final version of this section, summarize the quantitative findings, identify the best-performing configuration, and explain whether transfer learning provides a meaningful advantage for this task. You should also state what Grad-CAM reveals about the interpretability of the chosen model and discuss the practical value of combining performance analysis with visual explanation.

\section{Future Work}
Future work can extend this study in several directions. First, the binary classification setting can be expanded to multi-label classification so that individual thoracic diseases are modeled directly. Second, additional architectures such as EfficientNet, vision transformers, or medical-image-specific models can be evaluated. Third, more rigorous explainability and uncertainty estimation methods can be added to complement Grad-CAM. Finally, external validation on other chest X-ray datasets would help assess whether the conclusions generalize beyond the current benchmark.
